\documentclass[11pt]{article}
\usepackage{amsmath,amssymb,amsthm}
\newtheorem{theorem}{Theorem}
\title{Five Row Queries from Two Reverse Heads}
\author{}
\date{November 23, 2026}
\begin{document}
\maketitle

\paragraph{Status and limitation.}
This is a finite positive construction using independent column numerator
and denominator parameters. Its denominators have strictly negative
slopes and are positive on the entire Boolean cube. It does not vanish
on the coordinate axes, and it does not yield an unbounded number of
rows per head. Consequently it proves neither the requested local
compiler nor an improved asymptotic bound for worst-case head complexity.
The proof below is complete as a written argument, using the ordinary
real inverse function theorem. It has not been formalized in Lean.

\begin{theorem}
For every integer $s\geq1$ and every sign table
$\sigma:\{1,\ldots,5\}\times\{1,\ldots,s\}\to\{-1,1\}$, there
are affine numerators $N_1,N_2$ and affine denominators $B_1,B_2$
on the $(5+s)$-dimensional Boolean cube such that
\[
 \sigma(i,j)\left(\frac{N_1(e_i,e_j)}{B_1(e_i,e_j)}
                 +\frac{N_2(e_i,e_j)}{B_2(e_i,e_j)}\right)>0.
\]
Both denominators are positive everywhere on the cube and have strictly
negative slopes in every coordinate. For fixed $s$, all row coefficients
can be fixed independently of $\sigma$; only column coefficients vary.
\end{theorem}
\begin{proof}
Set $\varepsilon=1/(100(s+1))$, so $0<\varepsilon\leq1/200$. Define
\[
 (c_1,c_2,c_3,c_4,c_5)=(5,6,7,9,8),\qquad
 (k_1,k_2,k_3,k_4,k_5)=(6,8,5,7,11),
\]
and set $d_i=1-c_i\varepsilon$, $e_i=1-k_i\varepsilon$. For four
column parameters $x=(u,v,a,b)$, put
\[
 F_i(x)=\frac{A_i+a}{d_i-u}+\frac{C_i+b}{e_i-v}.
\]
The two distinguished parameter points are
\[
 x_+=(\varepsilon,2\varepsilon,3,5),\qquad
 x_-=(2\varepsilon,\varepsilon,5,3).
\]
For the first four rows take
\[
\begin{array}{c|cc}
 i&A_i&C_i\\\hline
 1&-3&-5\\
 2&-\dfrac{37\varepsilon-4}{17\varepsilon-2}
  &-3\dfrac{35\varepsilon-4}{17\varepsilon-2}\\
 3&-3\dfrac{31\varepsilon-4}{15\varepsilon-2}
  &-\dfrac{33\varepsilon-4}{15\varepsilon-2}\\
 4&-3\dfrac{39\varepsilon-4}{19\varepsilon-2}
  &-\dfrac{41\varepsilon-4}{19\varepsilon-2}.
\end{array}
\]
Direct substitution gives $F_i(x_+)=F_i(x_-)=0$ for $1\leq i\leq4$.
For the fifth row take
\[
 A_5=-\frac{2250\varepsilon^3-625\varepsilon^2+61\varepsilon-2}
                   {2\varepsilon(11\varepsilon-1)},\qquad
 C_5=\frac{2964\varepsilon^3-949\varepsilon^2+83\varepsilon-2}
                   {2\varepsilon(11\varepsilon-1)}.
\]
Another direct substitution gives $F_5(x_+)=1$ and $F_5(x_-)=-1$.

For $i\leq4$ clear the two denominators:
\[
 G_i(x)=(A_i+a)(e_i-v)+(C_i+b)(d_i-u).
\]
These are polynomials in the four column parameters. Their Jacobian
rows, with the parameter order $(u,v,a,b)$, are
\[
 (-C_i-b,\,-A_i-a,\,e_i-v,\,d_i-u).
\]
The determinants of the resulting four by four matrices at $x_+$ and
$x_-$ are, respectively,
\[
 \Delta=\frac{96\varepsilon^3
                 (337\varepsilon^2-74\varepsilon+4)}
              {(15\varepsilon-2)(17\varepsilon-2)(19\varepsilon-2)}
 \quad\text{and}\quad-\Delta.
\]
These identities follow by substituting the displayed row constants
and expanding the four by four determinants. They also provide short
exact-arithmetic verification certificates. None of their factors
vanishes: the three linear factors are negative, and
$337\varepsilon^2-74\varepsilon+4\geq4-74/200>0$.

All $d_i-u$ and $e_i-v$ are positive at both distinguished points.
For each choice of $\tau\in\{-1,1\}$, choose an open neighborhood
$U_\tau$ of $x_\tau$ small enough that these denominators stay positive,
$\tau F_5>0$, and
\[
 \varepsilon/2<u<5\varepsilon/2,\qquad
 \varepsilon/2<v<5\varepsilon/2.
\]
The inverse function theorem, shrinking this neighborhood if necessary,
shows that $G=(G_1,\ldots,G_4)$ maps a neighborhood of $x_\tau$
onto a neighborhood of zero in $\mathbb R^4$. For every prescribed
four signs, that neighborhood contains a sufficiently small vector
having those signs. Its inverse image has those signs for the first
four $F_i$, because the cleared denominators are positive; the fifth
sign is $\tau$. Hence all $32$ strict sign vectors occur with $u,v$
in the displayed bounds.

For each column $j$, choose a parameter vector
$(u_j,v_j,a_j,b_j)$ producing its desired five signs. Define on the
whole cube
\[
\begin{split}
 N_1(z,w)&=\sum_{i=1}^5 A_i z_i+\sum_{j=1}^s a_jw_j,
 &B_1(z,w)&=1-\varepsilon\sum_{i=1}^5c_i z_i-\sum_{j=1}^su_jw_j,\\
 N_2(z,w)&=\sum_{i=1}^5 C_i z_i+\sum_{j=1}^s b_jw_j,
 &B_2(z,w)&=1-\varepsilon\sum_{i=1}^5k_i z_i-\sum_{j=1}^sv_jw_j.
\end{split}
\]
Every slope of either denominator is strictly negative. Since
$\sum_i c_i=35$ and $\sum_i k_i=37$, the minimum of each denominator
on the whole cube is strictly greater than
\[
 1-37\varepsilon-\frac52s\varepsilon
 =1-\frac{37+(5/2)s}{100(s+1)}>0.
\]
On $(e_i,e_j)$ their sum of ratios is exactly
$F_i(u_j,v_j,a_j,b_j)$, giving the required signs.
\end{proof}

\paragraph{What would be needed for amplification.}
The two neighborhoods provide two regular parameter choices, while four
independently varying local coordinates supply four additional signs.
Repeating this finite gadget on separate row groups would require a
mechanism preventing the groups from interfering; none is supplied here.
Nor is an axis-zero transformation with a head saving supplied. The
ratio of five rows to two heads is fixed and is not an asymptotic
compression theorem.

\end{document}
